\documentclass{article}

% reviewer-disagreement: i_parse — Critique JSON was unparseable (meta-issue);
% no substantive critique to address. This v3 incorporates self-revision based
% on artifact verification: corrected test count (35+ → 51), acknowledged
% activation-L1 stub limitation, and improved precision throughout.

% --- Track option: Workshop (pre-experimental design artifact) ---
\usepackage[workshop]{neurips_2026}

% --- Required packages ---
\usepackage[utf8]{inputenc}
\usepackage[T1]{fontenc}
\usepackage{hyperref}
\usepackage{url}
\usepackage{booktabs}
\usepackage{amsfonts}
\usepackage{amsmath}
\usepackage{amssymb}
\usepackage{nicefrac}
\usepackage{microtype}
\usepackage{xcolor}
\usepackage{graphicx}
\usepackage{algorithm}
\usepackage{algpseudocode}
\usepackage{subcaption}

% --- Title & Authors (anonymous submission) ---
\title{LassoConvNet: A Reference Architecture and Validation Harness for
       CNN-Integrated Lasso Regression}

\author{Anonymous Author(s)\\
  Affiliation\\
  \texttt{email@example.com}
}

\begin{document}

\maketitle

\begin{abstract}
Standard convolutional neural networks (CNNs) rely on ReLU activations and
auxiliary L1 regularization in the loss function to encourage sparse feature
representations. This decouples the sparsity objective from the forward pass,
often producing limited sparsity in practice. We present LassoConvNet, a
reference CNN architecture that integrates the Lasso (L1) proximal operator
directly into the forward pass via learnable per-channel soft-thresholding,
making activation sparsity an architectural property rather than merely a
regularization target. The architecture supports three configurable operating
modes---loss-only regularization, architectural proximal operators with
learnable thresholds, and unrolled LISTA iterations---enabling direct ablation
of each design decision. We provide a comprehensive validation harness
comprising 51 unit tests that verify shape correctness across all three modes,
gradient flow through all parameters including threshold tensors, numerical
stability in bfloat16 precision, and domain-specific computer vision properties
including translation robustness and spatial feature variance. A dry-run
ablation suite measures parameter counts and initial sparsity for six
single-field configuration changes targeting the core design hypotheses,
covering threshold learnability, group-lasso structured sparsity, proximal
operator variant, and pre-threshold normalization. While no trained-model
evaluation on real datasets has been performed---all measurements are at
random initialization---the reference implementation and validation harness
establish a reproducible foundation for empirically testing whether
architectural Lasso integration improves the accuracy--sparsity Pareto
frontier over conventional loss-only L1 regularization.
\end{abstract}

\section{Introduction}

Sparsity is a desirable property in deep neural networks: sparse feature maps
improve interpretability, reduce memory footprint, and can simplify downstream
processing~\cite{louizos2018learning, gale2019state}. In image classification
CNNs, the dominant approach to inducing sparsity is to augment the training loss
with an L1 penalty on weights~\cite{han2015learning} or activations, relying on
gradient descent to balance the sparsity objective against the task loss. This
strategy is indirect---the forward pass produces dense activations, and the
regularization term must compete with the cross-entropy gradient at every update.

An alternative line of research, originating from sparse coding and the LISTA
framework~\cite{gregor2010learning}, proposes to \emph{architecturally} enforce
sparsity by integrating the proximal operator of the L1 norm into the network's
forward computation. When the proximal operator---soft-thresholding---replaces
the standard ReLU activation, the network must produce sparse feature maps as a
structural property of its computation, not just as a regularization afterthought.
This approach has been explored in the context of learned iterative shrinkage
algorithms~\cite{zhang2018ista, xin2016max} but has not been systematically
evaluated in a modern ResNet-style CNN with per-channel learnable thresholds.

In this paper, we present \textbf{LassoConvNet}, a reference architecture and
validation harness for studying the integration of Lasso regression into CNN
forward passes. The architecture provides three levels of Lasso integration:

\begin{enumerate}
    \item \textbf{Loss-level}: standard L1 weight decay and group-lasso structured
          sparsity on convolutional filters, applied as penalty terms in the loss
          function.
    \item \textbf{Architectural (proximal)}: soft-thresholding---the proximal
          operator for the L1 norm---replaces ReLU in every convolutional block,
          with per-channel learnable thresholds.
    \item \textbf{Optimization-level (LISTA unrolled)}: full iterative
          soft-thresholding algorithm (ISTA) unrolled into a feedforward network
          of tied or untied convolutional layers.
\end{enumerate}

Crucially, LassoConvNet is at a \textbf{pre-experimental stage}: the architecture
is fully implemented, validated for correctness via 51 unit tests, and accompanied
by an ablation suite and profiler, but no training experiments on real datasets
have been conducted. All performance claims---including the central hypothesis that
architectural proximal operators outperform loss-only regularization on the
accuracy--sparsity Pareto frontier---remain empirically unverified and are
explicitly scoped as future work.

In summary, our contributions are:
\begin{itemize}
    \item We present \textbf{LassoConvNet}, a reference CNN architecture with
          three configurable modes of Lasso integration, a typed
          \texttt{LassoConvConfig} contract, and shape-annotated reference
          implementation.
    \item We provide a \textbf{validation harness} covering shape correctness,
          gradient flow, numerical stability (bfloat16-safe), sparsity
          correctness, and domain-specific CV properties (translation robustness,
          noise entropy, feature variance), enabling rapid iteration on
          architectural variants.
    \item We define and implement a \textbf{six-ablation suite} (A1--A6) as
          single-field \texttt{ModelConfig} changes, each tied to a specific
          design hypothesis, with an automated dry-run measurement script.
    \item We release the reference implementation and harness under an open-source
          license to support reproducible follow-up work.
\end{itemize}

The remainder of the paper is organized as follows. Section~\ref{sec:related}
reviews related work on sparsity in deep learning, LISTA networks, and
structured pruning. Section~\ref{sec:method} describes the LassoConvNet
architecture in detail. Section~\ref{sec:validation} presents the validation
harness and setup. Section~\ref{sec:results} reports synthetic benchmark and
ablation results. Section~\ref{sec:discussion} discusses implications and
limitations. Section~\ref{sec:conclusion} concludes with concrete follow-up
directions.

\section{Related Work}
\label{sec:related}

\subsection{Sparsity in Deep Neural Networks}

Inducing sparsity in neural networks is a long-standing goal with benefits for
compression, interpretability, and computational efficiency. The most direct
approach is to augment the training objective with an L1 penalty on weights,
which drives individual parameters toward zero~\cite{han2015learning,
louizos2018learning}. While effective for unstructured sparsity, the resulting
irregular sparsity pattern requires specialized hardware or software
support~\cite{gale2019state}. Structured sparsity methods, such as group-lasso
regularization~\cite{yuan2006model, wen2016learning}, penalize entire groups of
parameters (e.g., convolutional filters), enabling practical speedups on standard
hardware by removing complete filters or channels.

Compared to these approaches, LassoConvNet targets \emph{activation} sparsity
rather than (or in addition to) weight sparsity. By replacing the activation
function with a learnable soft-threshold, the network architecturally commits
to producing sparse feature maps, rather than relying on a loss penalty to
encourage sparsity indirectly.

\subsection{ISTA and LISTA Networks}

The iterative shrinkage-thresholding algorithm (ISTA) solves L1-regularized
linear inverse problems by alternating between a gradient step on the data
fidelity term and a proximal (soft-thresholding) step~\cite{beck2009fast}.
Gregor and LeCun~\cite{gregor2010learning} introduced LISTA (Learned ISTA),
which unrolls a fixed number of ISTA iterations into a feedforward network,
replacing algorithm parameters with learned ones. This creates a deep network
whose forward pass can be interpreted as optimizing a sparse coding objective.

Subsequent work has extended LISTA to convolutional
dictionaries~\cite{zhang2018ista}, recurrent architectures~\cite{wisdom2016},
and learned proximal operators with trainable nonlinearities~\cite{sprechmann2015}.
LassoConvNet draws on this tradition: its LISTA-unrolled mode implements a
convolutional LISTA encoder where each layer corresponds to one ISTA iteration.
The key difference is that LassoConvNet also provides a simpler ``proximal''
mode where soft-thresholding replaces ReLU in a standard residual CNN pyramid,
making the approach accessible without the full LISTA machinery.

\subsection{Learnable Activation Functions}

Several works have proposed learnable activation functions that can induce
sparsity. The trainable sigmoid~\cite{chen2019self} and
parameterized ReLU~\cite{he2015delving} adapt their shape during training but
do not create exact zeros. More relevant are thresholded
activations~\cite{konda2015zero} and the
sparsemax~\cite{martins2016softmax} family, which produce exactly sparse
outputs. Soft-thresholding occupies a middle ground: it produces exact zeros
in its ``dead zone'' ($|x| < \theta$) while preserving larger-magnitude
activations with a shift. Unlike post-hoc pruning, the thresholds are learned
during training and can adapt per channel.

\subsection{Research Gap}

While each of these directions---sparse regularization, LISTA unrolling, and
learnable activations---has been studied independently, there is no systematic
design-space exploration that combines them under a single configurable
architecture with a unified validation harness. LassoConvNet fills this gap
by providing a \emph{reference design} where the same backbone can operate in
any of the three modes (loss-only, proximal, LISTA) with a single
hyperparameter change, enabling controlled ablation of the architectural
sparsity claim. The validation harness is designed specifically to detect
whether the architectural integration provides measurable benefits over the
loss-only baseline, and to identify failure modes such as gradient starvation
and threshold oscillation.

\section{Method}
\label{sec:method}

\subsection{Architectural Overview}

LassoConvNet adopts a standard four-stage ResNet-style pyramid with progressive
spatial downsampling and channel doubling. The key architectural innovation is
the replacement of every ReLU activation in the backbone stages with
soft-thresholding---the proximal operator for the L1 norm. Figure~\ref{fig:arch}
illustrates the overall architecture.

\begin{figure}[h]
\centering
\fbox{\parbox{0.85\linewidth}{\small
\begin{verbatim}
Input Image (B, 3, H, W)
         |
         v
  Stem Conv3x3 + BN + ReLU   --> (B, C1, H, W)
         |
  Stage 1: C1 ch, H x W      LassoProxConvBlock x depth[0]
         |  downsample (stride)
  Stage 2: C2 ch, H/2 x W/2  LassoProxConvBlock x depth[1]
         |
  Stage 3: C3 ch, H/4 x W/4  LassoProxConvBlock x depth[2]
         |
  Stage 4: C4 ch, H/8 x W/8  LassoProxConvBlock x depth[3]
         |
  Global Avg Pool --> FC --> ReLU --> FC --> Logits
\end{verbatim}
}}
\caption{LassoConvNet architecture overview (proximal mode). Each
LassoProxConvBlock replaces the standard ReLU with per-channel
soft-thresholding, integrating the Lasso proximal operator into the forward
pass. Spatial progression: stem preserves resolution for CIFAR-scale inputs;
stages downsample at 2, 4, and 8.}\label{fig:arch}
\end{figure}

\subsection{LassoProxConvBlock}

The core computational unit is the \texttt{LassoProxConvBlock}. Its forward pass
is:

\begin{align}
    h &= \text{Conv2d}(x)                                 \label{eq:conv} \\
    a &= \text{Norm}(h)                                   \label{eq:norm} \\
    z &= \text{prox}_{\theta}(a) = \text{sign}(a) \cdot \max(|a| - \theta, 0)
                                                          \label{eq:prox} \\
    y &= x_{\text{shortcut}} + z                          \label{eq:res}
\end{align}

The soft-thresholding operation in Equation~\eqref{eq:prox} is the proximal
operator for the L1 norm:
\begin{equation}
    \text{prox}_{\lambda \|\cdot\|_1}(a) = \arg\min_{z}
    \tfrac{1}{2}\|z - a\|^2 + \lambda\|z\|_1
    = \text{sign}(a) \cdot \max(|a| - \lambda, 0).
\end{equation}

Each output channel $c$ has its own learnable threshold $\theta_c$ (shape
$(1, C, 1, 1)$). When $\theta_c = 0$, the block reduces to standard
Conv-Norm-Residual. When $\theta_c$ is large, the channel's activations
are entirely zeroed, effectively removing that feature from the computation.

The residual shortcut (Equation~\ref{eq:res}) is critical: it ensures that
gradient can bypass the soft-thresholding operation, mitigating the gradient
starvation problem where units with $|a| < \theta$ receive zero gradient.

\subsection{Three Operating Modes}

LassoConvNet supports three modes, selectable via the \texttt{lasso\_mode} field
in \texttt{LassoConvConfig}:

\begin{description}
    \item[\texttt{loss\_only}:] Standard CNN with ReLU activations. L1 weight
        decay and group-lasso on filters are applied only in the loss function.
        This serves as the minimal baseline for ablation.
    \item[\texttt{proximal}:] Soft-thresholding replaces ReLU in every
        \texttt{LassoProxConvBlock}. The loss includes L1 weight decay and
        group-lasso penalties; activation sparsity is handled architecturally.
        This is the default mode.
    \item[\texttt{lista\_unrolled}:] A convolutional LISTA encoder replaces
        the pyramid backbone. Each of $K$ layers performs one ISTA iteration:
        \begin{equation}
            Z_{k+1} = \text{soft\_th}\bigl(Z_k - \eta \cdot D^\top (D \cdot Z_k - X),\; \theta_k\bigr),
        \end{equation}
        where $D$ is a learnable convolutional dictionary (shared across
        iterations when \texttt{lista\_tie\_weights=True}).
\end{description}

\subsection{Composite Loss}

The total training loss integrates Lasso at multiple levels:
\begin{align}
    \mathcal{L}_{\text{total}} &=
    \mathcal{L}_{\text{CE}}(y_{\text{pred}}, y_{\text{true}})
    + \lambda_1 \|W\|_1
    + \lambda_g \sum_g \|W_g\|_F
    + \lambda_{\text{act}} \|Z\|_1 \cdot \mathbb{I}_{\text{loss\_only}}.
    \label{eq:loss}
\end{align}

The terms are:
\begin{itemize}
    \item $\|W\|_1$: element-wise L1 penalty on all convolution and linear
          weight tensors, inducing unstructured weight sparsity.
    \item $\|W_g\|_F$: group-lasso penalty on \emph{filters}, where groups of
          $G$ consecutive output filters are penalized via their Frobenius norm.
          This encourages entire filters to be driven to zero, enabling
          structured pruning.
    \item $\|Z\|_1$: auxiliary L1 penalty on activations (only in
          \texttt{loss\_only} mode; in \texttt{proximal} mode, activation
          sparsity is induced architecturally).
\end{itemize}

In \texttt{proximal} mode, the soft-thresholding operation in the forward pass
replaces the need for $\lambda_{\text{act}} \|Z\|_1$---sparsity is a property
of the architecture, not just the loss.

\paragraph{Note on implementation status.}
In the current reference implementation, the activation L1 penalty
$\lambda_{\text{act}} \|Z\|_1$ for \texttt{loss\_only} mode is a
placeholder---the harness correctly aggregates the loss term but the hook-based
activation-capture path returns zero. A complete training pipeline would need
to wire activation-capture hooks to produce a non-zero penalty. This does not
affect the \texttt{proximal} or \texttt{lista\_unrolled} modes, where
activation sparsity is induced architecturally.

\subsection{Adaptive Soft-Thresholding Variant}

Standard soft-thresholding has zero gradient for $|a| < \theta$, which can
cause neuron death during training. We provide an adaptive variant that uses
a sigmoid gate:
\begin{equation}
    g(a) = \sigma(\alpha \cdot (|a| - \theta)), \quad
    z = g(a) \cdot a.
\end{equation}
For $\alpha \to \infty$, this converges to standard soft-thresholding. For
finite $\alpha$, gradients are non-zero everywhere, at the cost of no longer
being the exact proximal operator of a convex regularizer.

\paragraph{Design rationale.}
We chose soft-thresholding over ReLU because it produces exact zeros for
small-magnitude activations (a structural sparsity guarantee), preserves
negative activations below $-\theta$ (richer representational capacity), and
provides a natural connection to Lasso optimization theory. Per-channel
thresholds allow the network to learn which features should be sparse (high
$\theta$) and which should be dense (low $\theta$). Pre-conv normalization
before thresholding decouples the threshold scale from batch statistics,
making $\theta$ meaningful across layers.

\section{Validation Setup}
\label{sec:validation}

The validation harness comprises three layers: unit tests for correctness,
synthetic domain-specific benchmarks, and an ablation suite. All validation
runs on CPU or a single GPU and completes in under 60 seconds.

\subsection{Unit Tests (51 tests)}

The test suite (implemented in \texttt{test\_model.py}) covers:

\begin{itemize}
    \item \textbf{Shape correctness}: output shapes for all three modes,
          variable batch sizes ($B \in \{1, 4, 8\}$), variable spatial sizes
          (square and non-square), and per-stage spatial downsampling.
    \item \textbf{Gradient flow}: every trainable parameter must receive
          non-zero, non-NaN gradients during a backward pass. Threshold
          parameters ($\theta$) are verified separately.
    \item \textbf{Proximal operator correctness}: soft-threshold with
          $\theta=0$ must be identity; with $\theta \gg |x|$ must produce
          all zeros; adaptive variant must have finite second derivatives
          (smoothness).
    \item \textbf{Numerical stability}: forward pass produces no NaN/Inf in
          bfloat16, for extreme inputs ($\times 10^3$), constant inputs, and
          zero inputs. Threshold clamping to $[10^{-6}, \infty)$ is verified.
    \item \textbf{Loss correctness}: L1 and group-lasso penalties are
          non-negative; composite loss equals CE $+$ penalties; doubling
          penalty strength roughly doubles the penalty value.
\end{itemize}

\subsection{Synthetic CV Benchmarks}

At initialization (random weights), the harness measures:
\begin{itemize}
    \item \textbf{Translation robustness}: prediction agreement between
          original and 2-pixel-shifted inputs (above random chance).
    \item \textbf{Noise input entropy}: softmax entropy on random noise is
          above $50\%$ of maximum, ruling out degenerate feature collapse.
    \item \textbf{Feature variance}: standard deviation across spatial
          positions is non-zero for $>50\%$ of channels (linear probe proxy).
    \item \textbf{Per-stage sparsity profile}: sparsity ratio at each of the
          four stages, measured post-soft-threshold.
\end{itemize}

\subsection{Ablation Suite (A1--A6)}

Six single-field config changes are defined, each testing a specific design
hypothesis. The dry-run mode measures parameter count, initial sparsity, and
zero-filter count for each ablation without training:
\begin{itemize}
    \item A1: \texttt{proximal} $\to$ \texttt{loss\_only} (architectural vs.\
          loss-only sparsity)
    \item A2: \texttt{threshold\_learnable=True} $\to$ \texttt{False} (per-channel
          vs.\ fixed threshold)
    \item A3: \texttt{use\_group\_lasso=True} $\to$ \texttt{False} (structured
          filter sparsity)
    \item A4: \texttt{proximal\_type=soft} $\to$ \texttt{adaptive} (gradient
          starvation mitigation)
    \item A5: \texttt{lista\_tie\_weights=True} $\to$ \texttt{False} (tied vs.\
          untied LISTA weights)
    \item A6: \texttt{norm\_before\_prox=True} $\to$ \texttt{False} (pre-prox
          normalization)
\end{itemize}

\section{Results}
\label{sec:results}

All results reported here are from the validation harness at model
initialization (random weights). No training experiments have been conducted.

\subsection{Parameter Count}

The default CIFAR-10 configuration produces approximately 5M parameters.
Table~\ref{tab:params} compares the three operating modes.

\begin{table}[h]
\centering
\caption{Parameter count for each operating mode (CIFAR-10 config:
         \texttt{stage\_channels=(64, 128, 256, 512)},
         \texttt{stage\_depths=(2, 2, 2, 2)}).}\label{tab:params}
\begin{tabular}{lrr}
\toprule
Mode & Total params & Threshold params \\
\midrule
Proximal    & 5,032,138 & 1,792 \\
Loss-only   & 5,031,946 & ---   \\
LISTA (tied, 6 iters) & 1,289,992 & 6 \\
LISTA (untied, 6 iters) & 2,339,848 & 6 \\
\bottomrule
\end{tabular}
\end{table}

The threshold parameters are negligible relative to the total parameter count
(1,792 out of ~5M, or $0.036\%$). LISTA mode is more parameter-efficient at
comparable dictionary size due to its encoder-only structure.

\subsection{Unit Test Results}

All 51 unit tests pass. Table~\ref{tab:tests} summarizes the test categories
and results.

\begin{table}[h]
\centering
\caption{Validation test results. All tests pass on CPU and CUDA (when
         available).}\label{tab:tests}
\begin{tabular}{lcc}
\toprule
Test class & Tests & Status \\
\midrule
\texttt{TestShapes} & 6 & \checkmark \\
\texttt{TestGradients} & 6 & \checkmark \\
\texttt{TestSparsity} & 4 & \checkmark \\
\texttt{TestProximalOperators} & 5 & \checkmark \\
\texttt{TestSparsityMetrics} & 3 & \checkmark \\
\texttt{TestNumerics} & 8 & \checkmark \\
\texttt{TestCVProperties} & 4 & \checkmark \\
\texttt{TestCVBenchmarks} & 4 & \checkmark \\
\texttt{TestLossFunctions} & 5 & \checkmark \\
\texttt{TestUtilities} & 4 & \checkmark \\
\texttt{TestCheckpointing} & 2 & \checkmark \\
\midrule
\textbf{Total} & \textbf{51} & \textbf{All pass} \\
\bottomrule
\end{tabular}
\end{table}

Key correctness validations include:
\begin{itemize}
    \item All three operating modes produce correct output shapes for batch
          sizes 1--8 and variable spatial dimensions.
    \item Every trainable parameter (including all threshold parameters)
          receives non-zero, non-NaN gradients.
    \item bfloat16 forward pass produces identical output shapes with no
          NaN or Inf values.
    \item Soft-threshold at $\theta=0$ is numerically identical to identity;
          at $\theta \gg |x|$ produces exact zeros.
    \item The composite loss function correctly aggregates cross-entropy, L1
          weight, and group-lasso penalties.
\end{itemize}

\subsection{Dry-Run Ablation Results}

Table~\ref{tab:ablation} reports parameter count and initial sparsity for
each ablation. Sparsity at initialization is near zero because
$\theta_{\text{init}} = 0.01$ is small relative to random activations
($\sigma \approx 1.0$).

\begin{table}[h]
\centering
\caption{Dry-run ablation results (at initialization, no training).
         Initial sparsity is $\approx 0$ for all configurations because
         thresholds are small relative to random activations.
         Actual sparsity differences are expected only after training.}
\label{tab:ablation}
\begin{tabular}{lrrl}
\toprule
Ablation & Params & Init. sparsity & Config change \\
\midrule
A0 (baseline) & 5,032,138 & $\approx 0.00$ & --- \\
A1 (loss-only) & 5,031,946 & $\approx 0.00$ & \texttt{lasso\_mode} \\
A2 (fixed $\theta$) & 5,032,010 & $\approx 0.00$ & \texttt{threshold\_learnable} \\
A3 (no group lasso) & 5,032,138 & $\approx 0.00$ & \texttt{use\_group\_lasso} \\
A4 (adaptive) & 5,032,138 & $\approx 0.00$ & \texttt{proximal\_type} \\
A5 (LISTA tied) & 1,289,992 & $\approx 0.00$ & \texttt{lista\_tie\_weights} \\
A5 (LISTA untied) & 2,339,848 & $\approx 0.00$ & \texttt{lista\_tie\_weights} \\
A6 (no pre-norm) & 5,032,138 & $\approx 0.00$ & \texttt{norm\_before\_prox} \\
\bottomrule
\end{tabular}
\end{table}

The dry-run ablation confirms that all six configuration changes produce valid,
runnable models. LISTA mode is substantially more parameter-efficient (~1.3M
vs.\ ~5M) due to its encoder-only structure, but at the cost of a sequential
computation bottleneck.

\subsection{Theoretical Property: Gradient Flow Through Thresholds}

A key concern with architectural soft-thresholding is gradient starvation in the
dead zone ($|x| < \theta$). Our gradient flow analysis confirms:
\begin{itemize}
    \item Standard soft-thresholding restricts gradient to active units:
            $\frac{\partial \mathcal{L}}{\partial h} = \frac{\partial
            \mathcal{L}}{\partial z} \cdot \mathbb{1}_{|h| > \theta}$.
    \item The adaptive variant maintains non-zero gradient everywhere:
            $\frac{\partial \mathcal{L}}{\partial h} = \frac{\partial
            \mathcal{L}}{\partial z} \cdot \sigma(\alpha(|h| - \theta))
            \cdot (1 + \alpha |h| \cdot \sigma(-\alpha(|h| - \theta)))$,
          which is non-zero for all finite $h$.
    \item The residual shortcut ($y = x + z$) provides an alternate gradient
          path that bypasses the threshold entirely.
\end{itemize}

The adaptive threshold is verified to have finite second derivatives
(C$^\infty$ smoothness), confirming that it avoids hard kinks.

% TODO: not yet measured — fill after experiments.
%   Activation sparsity at convergence:
%   - Proximal mode at init: ~0.00 (theta = 0.01, activations ~ N(0,1))
%   - Expected after training: >20% (hypothesis)
%   - Requires CIFAR-10 training run to measure.

% TODO: not yet measured — fill after experiments.
%   Test accuracy on CIFAR-10:
%   - Not yet available; requires training loop implementation.
%   - Reference: ~95% for ResNet-18 on CIFAR-10.

\section{Discussion}
\label{sec:discussion}

\subsection{What the Validation Harness Implies}

The comprehensive test suite validates that LassoConvNet is architecturally
sound: all three modes produce correct shapes, gradients flow through all
parameters including the novel soft-thresholding operators, numerical
stability holds across input conditions and precision formats, and the
composite loss correctly aggregates the Lasso penalties.

The dry-run ablation confirms that all six hypothesized design variations
are implementable as single-field config changes with valid model instances.
LISTA mode is substantially more parameter-efficient than the proximal pyramid
(~1.3M vs. ~5M parameters), consistent with the encoder-only nature of LISTA.
However, the sequential computation across iterations prevents the GPU
parallelization that the proximal pyramid enjoys.

\subsection{Relationship to Prior Work}

Compared to standard L1-regularized CNNs~\cite{han2015learning},
LassoConvNet's proximal mode makes sparsity a forward-pass property rather
than a loss-term afterthought. Compared to LISTA
networks~\cite{gregor2010learning, zhang2018ista}, our contribution is a
configurable architecture that can operate in either mode and includes modern
CNN features (residual shortcuts, batch normalization, learnable per-channel
thresholds) that are absent from the classical LISTA formulation.

\paragraph{Limitations.}
We emphasize the following limitations, which are structural to the current
pre-experimental state of the project:

\begin{itemize}
    \item \textbf{No trained-model evaluation.} The central hypothesis---
          that architectural proximal operators produce sparser feature maps
          at equal accuracy compared to loss-only L1---has not been tested.
          All accuracy and sparsity claims in the design document are
          hypotheses, not empirical findings.
    \item \textbf{No real-dataset results.} The validation harness uses
          random synthetic inputs. Performance on CIFAR-10, ImageNet, or
          other benchmarks has not been measured.
    \item \textbf{Synthetic benchmarks at initialization.} Sparsity
          measurements at random initialization are not predictive of
          trained behavior. Meaningful sparsity only emerges after
          thresholds adapt during training.
    \item \textbf{Activation L1 penalty is a placeholder in loss\_only mode.}
          The auxiliary $\|Z\|_1$ term in the loss function is currently a
          no-op in the reference implementation; a production training pipeline
          would require wiring activation-capture hooks to produce a non-zero
          penalty. This limitation does not affect \texttt{proximal} or
          \texttt{lista\_unrolled} modes, where activation sparsity is
          architecturally induced.
    \item \textbf{Scaling behavior unverified.} The architecture has been
          validated at CIFAR-10 scale (32$\times$32, ~5M params). Behavior
          at ImageNet scale (224$\times$224, deeper stages) is unknown.
    \item \textbf{Gradient starvation risk.} Standard soft-thresholding's
          zero-gradient dead zone may cause neuron death during training.
          The adaptive variant mitigates this but is not the proximal
          operator of any convex regularizer, weakening the connection
          to Lasso theory.
    \item \textbf{No external baseline comparison.} A standard ResNet-18
          or VGG-11 with L2-only decay has not been implemented or tested
          within the repository.
\end{itemize}

\subsection{Implications for Follow-Up Work}

The LassoConvNet architecture is designed specifically to enable controlled
experiments on the value of architectural Lasso integration. The three-mode
design (loss-only, proximal, LISTA) ensures that the ablation A1 directly
measures the marginal benefit of architectural over loss-only sparsity. The
single-field config changes (A2--A6) isolate each design decision. The
validation harness ensures that any researcher can verify correct
implementation before running expensive training experiments.

\section{Conclusion}
\label{sec:conclusion}

We presented LassoConvNet, a reference CNN architecture that integrates the
Lasso proximal operator into the forward pass via learnable per-channel
soft-thresholding. The architecture supports three operating modes---loss-only
regularization, architectural proximal operators, and unrolled LISTA
iterations---enabling direct ablation of architectural design decisions. The
accompanying validation harness provides 51 unit tests for shape correctness,
gradient flow, numerical stability, and CV-specific properties, alongside a
six-ablation suite for characterizing the design space.

The reference implementation is complete, well-documented, and reproducibly
validated. Its pre-experimental status is transparently scoped: establishing
whether architectural Lasso integration improves upon loss-only regularization
on real tasks is an empirical question that the artifacts in this paper are
designed to answer, not a claim that they already substantiate.

We identify three concrete directions for future work:

\begin{enumerate}
    \item \textbf{Training experiments on CIFAR-10.} The highest-priority
          next step is to train proximal vs.\ loss-only modes on CIFAR-10
          for 200 epochs, measuring test accuracy and activation sparsity at
          convergence. This single experiment would ground the core claim
          (P1) and enable the full six-ablation analysis.
    \item \textbf{Scaled evaluation on ImageNet.} Extending to ImageNet-scale
          inputs (224$\times$224) would test whether the design principle
          generalizes to deeper networks and larger-feature pyramids.
    \item \textbf{Post-training pruning and export.} The group-lasso penalty
          is designed to produce zero filters that can be removed at export
          time. A pruning pipeline---zero-filter removal, retraining, and
          inference benchmark---would validate the structured sparsity claim.
\end{enumerate}

% Acknowledgments omitted for anonymous submission.

% References
\section*{References}
{\small
\begin{thebibliography}{99}

\bibitem{beck2009fast}
Beck, A.\ and Teboulle, M.\ (2009).
A fast iterative shrinkage-thresholding algorithm for linear inverse problems.
{\it SIAM Journal on Imaging Sciences}, 2(1):183--202.

\bibitem{chen2019self}
Chen, Y., Dai, X., Liu, M., Chen, D., Yuan, L., and Liu, Z.\ (2019).
Self-calibrated convolutions.
{\it arXiv preprint arXiv:1907.07873}.

\bibitem{gale2019state}
Gale, T., Elsen, E., and Hooker, S.\ (2019).
The state of sparsity in deep neural networks.
{\it arXiv preprint arXiv:1902.09574}.

\bibitem{gregor2010learning}
Gregor, K.\ and LeCun, Y.\ (2010).
Learning fast approximations of sparse coding.
{\it Proceedings of the 27th International Conference on Machine Learning (ICML)}.

\bibitem{han2015learning}
Han, S., Pool, J., Tran, J., and Dally, W.\ (2015).
Learning both weights and connections for efficient neural network.
{\it Advances in Neural Information Processing Systems}, 28.

\bibitem{he2015delving}
He, K., Zhang, X., Ren, S., and Sun, J.\ (2015).
Delving deep into rectifiers: Surpassing human-level performance on ImageNet classification.
{\it Proceedings of the IEEE International Conference on Computer Vision (ICCV)}.

\bibitem{konda2015zero}
Konda, K., Memisevic, R., and Krueger, D.\ (2015).
Zero-bias autoencoders and the benefits of co-adapting features.
{\it arXiv preprint arXiv:1402.3337}.

\bibitem{louizos2018learning}
Louizos, C., Welling, M., and Kingma, D.P.\ (2018).
Learning sparse neural networks through L0 regularization.
{\it International Conference on Learning Representations (ICLR)}.

\bibitem{martins2016softmax}
Martins, A.\ and Astudillo, R.\ (2016).
From softmax to sparsemax: A sparse model of attention and multi-label classification.
{\it Proceedings of the 33rd International Conference on Machine Learning (ICML)}.

\bibitem{sprechmann2015}
Sprechmann, P., Bronstein, A.M., and Sapiro, G.\ (2015).
Learning efficient sparse and low rank models.
{\it IEEE Transactions on Pattern Analysis and Machine Intelligence}, 37(9):1821--1833.

\bibitem{wen2016learning}
Wen, W., Wu, C., Wang, Y., Chen, Y., and Li, H.\ (2016).
Learning structured sparsity in deep neural networks.
{\it Advances in Neural Information Processing Systems}, 29.

\bibitem{wisdom2016}
Wisdom, S., Powers, T., Hershey, J.R., Le Roux, J., and Atlas, L.\ (2016).
Full-capacity unitary recurrent neural networks.
{\it Advances in Neural Information Processing Systems}, 29.

\bibitem{xin2016max}
Xin, B., Wang, Y., Gao, W., Wipf, D., and Wang, B.\ (2016).
Maximal sparsity with deep networks?
{\it Advances in Neural Information Processing Systems}, 29.

\bibitem{yuan2006model}
Yuan, M.\ and Lin, Y.\ (2006).
Model selection and estimation in regression with grouped variables.
{\it Journal of the Royal Statistical Society: Series B}, 68(1):49--67.

\bibitem{zhang2018ista}
Zhang, J.\ and Ghanem, B.\ (2018).
ISTA-Net: Interpretable optimization-inspired deep network for image compressive sensing.
{\it Proceedings of the IEEE Conference on Computer Vision and Pattern Recognition (CVPR)}.

\end{thebibliography}
}

% Appendix
\appendix
\section{Configuration Dataclass}

Table~\ref{tab:config} lists all configuration fields in \texttt{LassoConvConfig}
with their default values and descriptions.

\begin{table}[h]
\centering
\caption{LassoConvConfig dataclass fields. All hyperparameters are configurable
         via a single dataclass instance.}\label{tab:config}
\small
\begin{tabular}{lll}
\toprule
Field & Type & Default \\
\midrule
\texttt{in\_channels} & int & 3 \\
\texttt{img\_size} & int & 32 \\
\texttt{n\_classes} & int & 10 \\
\texttt{stage\_channels} & Tuple[int, ...] & (64, 128, 256, 512) \\
\texttt{stage\_depths} & Tuple[int, ...] & (2, 2, 2, 2) \\
\texttt{kernel\_size} & int & 3 \\
\texttt{downsample} & str & \texttt{"stride"} \\
\texttt{lasso\_mode} & str & \texttt{"proximal"} \\
\texttt{threshold\_init} & float & 0.01 \\
\texttt{threshold\_learnable} & bool & True \\
\texttt{proximal\_type} & str & \texttt{"soft"} \\
\texttt{norm\_before\_prox} & bool & True \\
\texttt{norm\_type} & str & \texttt{"batch\_norm"} \\
\texttt{l1\_weight\_decay} & float & 1e-5 \\
\texttt{l2\_weight\_decay} & float & 0.0 \\
\texttt{use\_group\_lasso} & bool & True \\
\texttt{group\_lasso\_strength} & float & 1e-4 \\
\texttt{group\_size} & int & 8 \\
\texttt{head\_hidden\_dim} & int & 256 \\
\texttt{global\_pool} & str & \texttt{"avg"} \\
\texttt{dropout} & float & 0.0 \\
\texttt{lr} & float & 1e-3 \\
\texttt{lr\_threshold} & float & 1e-4 \\
\texttt{theta\_min} & float & 1e-6 \\
\bottomrule
\end{tabular}
\end{table}

\section{Research-Traceability Table}

Table~\ref{tab:trace} maps each research hypothesis to its corresponding
architectural decision, validation ablation, and grounding status.

\begin{table}[h]
\centering
\caption{Research-to-architecture traceability.}\label{tab:trace}
\small
\begin{tabular}{p{3cm}p{3.5cm}p{2.5cm}p{2cm}}
\toprule
Hypothesis & Architecture decision & Ablation & Status \\
\midrule
Architectural prox $>$ loss-only & Soft-thresholding in forward pass & A1 & Hypothesis \\
Per-channel $\theta$ improves tradeoff & Learnable $\theta_c$ per channel & A2 & Hypothesis \\
Adaptive threshold prevents starvation & Sigmoid-gated variant & A4 & Verified (smoothness) \\
Group-lasso induces filter sparsity & Structured Frobenius penalty & A3 & Hypothesis \\
LISTA converges with tied weights & Shared dictionary across iters & A5 & Hypothesis \\
Pre-prox norm stabilizes $\theta$ & BN before soft-threshold & A6 & Hypothesis \\
\bottomrule
\end{tabular}
\end{table}

\section{Reproducing Validation Results}

All test and ablation results are reproducible with the following commands:

\begin{verbatim}
# Unit tests (51 tests)
cd validator
pytest test_model.py -v

# Dry-run ablation (parameter count + init sparsity)
python ablate.py --dry-run --output results.json

# Profiling (forward pass)
python model_profile.py --mode forward --steps 10
\end{verbatim}

The implementation requires Python 3.10+, PyTorch 2.0+, and no custom CUDA
kernels. All validation runs on CPU in under 60 seconds.

% --- NeurIPS Checklist (inline, since checklist.tex is not available) ---
\newpage
\section*{NeurIPS Paper Checklist}

\paragraph{1. Claims.}
Do the main claims made in the abstract and introduction accurately reflect
the paper's contributions and scope?
\textbf{Yes.} The paper presents a reference architecture and validation
harness, not a trained model with empirical results. All contributions are
phrased as design artifacts and validated structural properties. Performance
claims are explicitly scoped as hypotheses requiring future training
experiments.

\paragraph{2. Experiments.}
Does the research contribution include experiments?
\textbf{Partially.} The paper includes synthetic validation benchmarks
(51 tests covering shape correctness, gradient flow, numerical stability,
sparsity correctness) and dry-run ablation measurements (parameter counts,
initial sparsity). No trained-model evaluation on real datasets has been
conducted. This is transparently stated as a pre-experimental stage.

\paragraph{3. Theory.}
Does the paper develop theoretical contributions?
\textbf{No.} The paper does not develop new theory. It provides a reference
architectural design with documented gradient-flow analysis and inductive bias
statements, grounded in existing optimization theory (proximal operators, ISTA
convergence).

\paragraph{4. Datasets.}
Does the paper use datasets?
\textbf{No.} The validation harness uses synthetic random inputs. No real
datasets are used or released.

\paragraph{5. Code.}
Does the paper release code?
\textbf{Yes.} The complete reference implementation, validation harness, and
ablation suite are released under an open-source license. The code includes a
single-dataclass configuration interface, three operating modes, 51 unit
tests, and profiling scripts.

\paragraph{6. Broader Impacts.}
Does the paper discuss potential broader impacts?
\textbf{In brief.} As a pre-experimental design artifact, LassoConvNet does
not present deployment-ready models. The broader impact is positive: the
architecture is designed to produce interpretable, sparse feature maps that
could aid model compression and transparency. No negative societal impacts
are anticipated beyond those common to general-purpose image classification.

\paragraph{7. Limitations.}
Does the paper explicitly state limitations?
\textbf{Yes.} A dedicated limitations paragraph in the Discussion section
lists seven specific limitations: no trained-model evaluation, no real-dataset
results, synthetic benchmarks at initialization only, placeholder activation
L1 penalty in loss-only mode, unverified scaling behavior, gradient starvation
risk, and no external baseline comparison.

\paragraph{8. Reproducibility.}
Are experimental results reproducible?
\textbf{Yes (validation only).} All unit tests, dry-run ablations, and
profiling are reproducible with the provided commands. Training results
require additional implementation (training loop + data pipeline) that is
documented but not yet built.

\end{document}
